% -----------------------------
% 執筆にあたっての留意点
% -----------------------------
%
% 関連研究の章では，この論文で提案・主張することの意義を「相対的」に位置づけるために，関連する研究を紹介します．決して「類似する研究」を並べるわけではありません．自分の論文をうまく位置づけるためにも，3つ程度の観点（節）を設けて，関連研究を紹介するのがよいでしょう．

% 例えば，「情報の信憑性を積極的にチェックするための動機付けを行う方法」に関する論文を書いているならば，関連研究を整理する観点として

% ・信憑性が怪しい情報を計算機がチェックする方法に関する研究
% ・人間が自分の力で信憑性をチェックできるようになることを支援する研究
% ・そもそも人間がどのように信憑性をチェックしているのかを明らかにする研究

% などが観点として考えられます．

% 各観点（節）では最低でも3報は関連する論文を引用し，それぞれの論文で何を提案しているのか，何が明らかにされたのかを1，2行程度でまとめてください．その上で，各節の最後で，それら論文と自分の論文の方向性の違いを主張してください．
%
%
% -----------------------------
% 補足（参考：how-to-xx/how-to-write-a-paper/1-学術論文の書き方.md 4章）
% -----------------------------
%
% ■ 各段落の主役は，先行研究ではなく「自分の研究との関係」です
% 「Aは〜した．Bは〜した．Cは〜した」と並べただけでは，この章の役割を果たしていません．
% 先行研究という地図の中に自分の研究を位置づけ，「既存研究では埋まっていない穴があり，本研究がそれを埋める」
% ことを示すのが目的です．各節では，そのグループの到達点と限界を述べた上で，節の最後に
% 「これらに対し本研究は〜の点で異なる」と差分を明示してください．
%
% ■ 先行研究を貶めてはいけません
% 先行研究への敬意は，あなたの論文の価値を下げません．むしろ比較対象が優れているほど，
% それを超えるあなたの貢献が際立ちます．また現実問題として，査読者がその先行研究の著者本人である
% 可能性は決して低くありません．「〜には限界がある」と書くのは構いませんが，
% 「〜は不十分である」「〜は誤っている」といった断定は，根拠を示せない限り避けてください．
%
% ■ 引用は，必ず自分で原典を読んでから行ってください
% 他の論文の関連研究の記述をそのまま孫引きすると，誤った要約をそのまま引き継ぐことになります．
% 少なくともアブストラクトと図，結論には目を通してから引用しましょう．
%
% ■ この章をどこに置くか
% 本テンプレートでは2章（イントロダクションの直後）に置いていますが，末尾（考察と結論の間）に置く流儀もあります．
% 判断基準は次の通りです．
%
% * 冒頭に置く：短く，かつ十分詳細に書ける場合．または，先行研究に対して早い段階で自分の立場を明確にしたい場合
% * 末尾に置く：イントロダクションで手短に触れておける場合．または，先行研究と十分に比較するために提案手法の技術的な説明が必要な場合
%
% 迷ったら，投稿先の学会の採録論文を数本見て，その分野の慣習に合わせるのが安全です．


% ----------------------
% 関連研究/Related Work
% ----------------------

\section{関連研究}
